\section{序列到序列模型与瓶颈问题}

\subsection{Encoder-Decoder RNN 回顾}

序列到序列（Seq2Seq）问题的经典架构是 \textbf{Encoder-Decoder RNN}：

\begin{figure}[H]
\centering
\includegraphics[width=0.95\textwidth]{slides-images/slide-003.jpg}
\caption{Encoder-Decoder RNN 架构：编码器 RNN 将输入序列压缩为上下文向量 $c$，解码器 RNN 基于 $c$ 生成输出序列\protect\footnotemark}
\end{figure}
\footnotetext{Slides 第4页。}

\begin{enumerate}
    \item \textbf{编码器}（Encoder）：一个 RNN，逐词处理输入序列（如英文句子``we see the sky''），将整个输入压缩为最终隐藏状态
    \item \textbf{上下文向量}（Context Vector）$c$：编码器的最终隐藏状态，被视为整个输入序列的``摘要''
    \item \textbf{解码器}（Decoder）：另一个 RNN，以上下文向量 $c$ 为初始输入，逐词生成输出序列（如意大利语翻译）
\end{enumerate}

\subsection{固定长度瓶颈}

\begin{figure}[H]
\centering
\includegraphics[width=0.95\textwidth]{slides-images/slide-004.jpg}
\caption{信息瓶颈问题：整个输入序列的信息必须压缩到一个固定长度的向量 $c$ 中\protect\footnotemark}
\end{figure}
\footnotetext{Slides 第5页。}

\begin{warningbox}{固定长度向量的信息瓶颈}
问题的核心在于：输入序列和输出序列之间的\textbf{唯一通信通道}是上下文向量 $c$——一个固定维度的向量（如 128 或 1024 个浮点数）。

\begin{itemize}
    \item 对于短句子（``we see the sky''），1024 维可能足够
    \item 对于长段落、整本书或海量语料，将所有信息压缩到 1024 个浮点数中显然不切实际
    \item 随着输入长度增加，信息损失会越来越严重
\end{itemize}
\end{warningbox}

\textbf{解决方案的直觉}：不再强制网络将整个输入序列压缩为一个向量。取而代之，让解码器在生成每个输出 token 时，都有机会\textbf{回头查看}整个输入序列的不同部分——这就是注意力机制的核心思想。

\subsection{本章小结}

Encoder-Decoder RNN 中的固定长度上下文向量是一个严重的信息瓶颈。注意力机制的诞生正是为了消除这个瓶颈，让解码器能够在每个时间步动态地访问编码器的所有隐藏状态。

%% ========== RNN 中的注意力 ==========

